\documentclass[11pt]{letter}

\usepackage[margin=0.5in]{geometry}
\usepackage{hyperref}

\begin{document}

\begin{letter}{}

\opening{Dear Arman,}

I am writing to apply for the Junior Full Stack Developer role at SayVo AI. With a Master of Science in Computer Science from the University of Calgary and hands-on experience building Python and JavaScript software, backend APIs, data pipelines, and LLM-enabled automation tools, I am excited by the chance to help build AI employees that call, qualify, book meetings, and connect reliably with real client workflows.

What stood out to me about SayVo is the ownership in the role. You are not looking for someone to sit behind a narrow ticket queue; you are looking for a builder who can write product code, connect systems together, debug messy integrations, explain technical details clearly, and keep moving when a small startup needs many things at once. That environment fits the way I like to work: I enjoy learning fast, tracing problems carefully, and turning vague requirements into working tools.

My backend experience started with Django REST APIs for NLP, text-to-speech, and speech-to-text services backed by PostgreSQL. I built authentication flows, payment-service logic, and plan-management features, and I also contributed to a WebRTC-based chatbot system using GPT-style responses and FAQ data. That work gave me a practical foundation in APIs, databases, user-facing backend logic, and the kind of debugging needed when different services and user flows need to behave consistently.

More recently, I have been building an agentic job-search automation platform using LLM workflows, n8n, OpenClaw, Docker, Google Sheets, Google Drive, SearXNG, and Telegram API integrations. The project involves prompt-driven analysis, workflow automation, structured outputs, spreadsheet tracking, and end-to-end automation around real decisions. While it is not a voice-agent product, it is closely aligned with SayVo's need for someone who can wire AI workflows together, think carefully about prompts, connect APIs, and improve a system based on real usage.

My research software work also strengthened my engineering discipline. As a research software developer with BigGeo, Mitacs, and the University of Calgary, I built Python and C++ pipelines for large-scale data ingestion, preprocessing, storage, retrieval, and performance evaluation. I improved one C++ processing workflow from 233 seconds to 26 seconds for a 10-tensor workload and reduced storage by 38\% on a scaled Canada-wide dataset. That experience taught me to care about correctness, performance, documentation, and small technical details that make systems reliable.

I have not yet worked directly with Twilio, A2P/10DLC, GoHighLevel, Pipedrive, Bonzo, or mortgage sales workflows, but I am comfortable learning new tools quickly and digging through documentation until I understand how a system really works. I also enjoy explaining technical ideas in plain English, which I know matters when onboarding clients and helping non-technical users trust an AI workflow.

Thank you for considering my application. I would be excited to bring my backend development, AI automation, prompt-workflow, and product-building experience to SayVo AI and grow with the company from this early stage.

\closing{Sincerely,

Amir Mirzai Golpayegani

\href{mailto:amirgolpaa24@gmail.com}{amirgolpaa24@gmail.com}}

\end{letter}

\end{document}
